\section{Conclusion}

We have developed a type theory for semilattice-enriched representable 
multicategories and shown that it has the standard desirable properties of a 
well-behaved type theory: terms uniquely determine their derivations, substitution 
is admissible and satisfies associativity and interchange, and the type theory 
correctly defines the free (representable) $\catname{SLat}$-multicategory over a 
multigraph.
The connection to e-graphs is made precise by showing that the $+$ operation of 
the type theory corresponds to the equivalence classes of e-graphs, and that 
equality saturation can be formulated as a fixed point of term rewriting modulo 
the structural equations of the type theory.

The type theory naturally accommodates extensions that capture richer 
computational phenomena present in e-graphs.
The explicit \textbf{let} construct captures sharing of subterms, corresponding 
to the node-sharing structure of e-graphs.
The recursive \textbf{letrec} construct captures cyclic e-graphs, which arise 
when equality saturation introduces cycles, and is modelled categorically by 
traced monoidal $\catname{SLat}$-categories.

We hope that this work contributes to a cleaner mathematical foundation for 
e-graphs and equality saturation, complementing the existing categorical accounts 
of~\citet{tiurin2025equivalencehypergraphsdporewriting} at the level of syntax 
rather than string diagrams.

\paragraph{Future work.}
On the computational side, it would be interesting to investigate whether the 
type theory can be used as the basis for an implementation of equality saturation 
that is more amenable to formal verification than existing algorithms.

Finally, we speculate that this work may lay some of the mathematical groundwork 
for programming languages with e-graphs built in as a primitive concept.
Just as the simply typed lambda calculus underlies the design of functional 
programming languages, a type theory of the kind developed here could inform the 
design of a language where equality saturation is not an external optimization 
pass but an intrinsic part of the language's evaluation semantics --- one where 
terms natively denote equivalence classes of programs rather than individual 
programs.
We leave the development of such a language as an ambitious direction for future 
work.